\section{Preliminaries}
\label{sec:prelim}

We collect the definitions and standard results used throughout. All random
variables are discrete with finite support, and all logarithms are base $2$, so
information is measured in bits. We write $H(\cdot)$ for Shannon entropy,
$H(\cdot\mid\cdot)$ for conditional entropy, $\II(\cdot\,;\cdot)$ for mutual
information and $\II(\cdot\,;\cdot\mid\cdot)$ for conditional mutual information, in
the conventions of \cite{cover2006elements}. We write $X\indep Y\mid Z$ for the
conditional independence of $X$ and $Y$ given $Z$.

\subsection{The information model}

\begin{definition}[Aligned behavior and pathways]
\label{def:behavior}
An \emph{aligned behavior} is a discrete random variable $S$ taking values in a finite
set $\mathcal{S}$ (for instance a refuse/comply decision, or a quantized
refusal logit-difference). A \emph{pathway readout} is a random variable
$R_i=\phi_i(\text{activations})$, a deterministic function of the model's internal
activations---an attention head, a projection onto a refusal direction, a group of
sparse-autoencoder latents, or an entire layer. We collect the readouts into the
vector $R=(R_1,\dots,R_n)$ and write $R_{-i}=(R_j)_{j\neq i}$ and, for $T\subseteq
\{1,\dots,n\}$, $R_T=(R_i)_{i\in T}$ and $T^c=\{1,\dots,n\}\setminus T$.
\end{definition}

\begin{definition}[Ablation and information cost]
\label{def:ablation}
To \emph{ablate} the set $T$ is to replace $\{R_i\}_{i\in T}$ by a fixed constant
value, so that the behavior must be read from the surviving readouts $R_{T^c}$. The
\emph{information cost} of ablating $T$ is
\[
  \Delta I_T \;\coloneqq\; \II(S;R) - \II(S;R_{T^c}),
\]
the drop in the information that the available readouts carry about $S$. We write
$\Delta I_i \coloneqq \Delta I_{\{i\}}$ for a single-pathway ablation.
\end{definition}

\begin{definition}[Conditional redundancy and $k$-redundancy]
\label{def:redundancy}
A pathway $R_i$ is \emph{redundant given the rest} if $S\indep R_i\mid R_{-i}$. The
behavior $S$ is \emph{$k$-redundant} if $\II(S;R_T)=\II(S;R)$ for every $T$ with
$\lvert T\rvert\ge n-k+1$; that is, every subset of at least $n-k+1$ pathways carries
all the information $R$ has about $S$.
\end{definition}

\begin{definition}[Behavioral fidelity]
\label{def:fidelity}
Let $\widehat{S}=g(R_{T^c})$ be a decoder of the behavior from the surviving
readouts, with error probability $P_e\coloneqq\Pr[\widehat{S}\neq S]$. For a
vector-valued behavior with pre- and post-ablation outputs $b_{\mathrm{pre}},
b_{\mathrm{post}}$, the \emph{output robustness} is
$\rho\coloneqq 1-\lVert b_{\mathrm{post}}-b_{\mathrm{pre}}\rVert_2/\lVert
b_{\mathrm{pre}}\rVert_2$, so that $\rho=1$ means behavior is perfectly preserved.
\end{definition}

\subsection{The graph model}

\begin{definition}[Computation DAG and pathway redundancy]
\label{def:dag}
A \emph{computation DAG} is a directed acyclic graph $G=(V,E)$ with a distinguished
source $\sigma$ and sink $\tau$. The behavior $S$ is \emph{computable} if at least one
directed $\sigma\to\tau$ path is intact. The \emph{pathway redundancy} of $G$ is its
$\sigma$--$\tau$ vertex connectivity,
\[
  R(G)\;\coloneqq\;\kappa_{\sigma\tau}(G),
\]
the minimum number of internal vertices whose deletion separates $\sigma$ from
$\tau$.
\end{definition}

\begin{definition}[The bundle circuit $B(w,d)$]
\label{def:bundle}
For integers $w,d\ge 1$, the \emph{bundle circuit} $B(w,d)$ consists of a source
$\sigma$, a sink $\tau$, and $d$ \emph{bundles} arranged in series, each bundle a set
of $w$ parallel internal vertices. Consecutive bundles are joined by a complete
bipartite set of edges, and $\sigma$ (resp.\ $\tau$) is joined to every vertex of the
first (resp.\ last) bundle. This is the canonical exactly solvable redundant circuit:
$w$ is its \emph{redundancy width} and $d$ its \emph{depth}.
\end{definition}

\begin{definition}[Random and targeted ablation thresholds]
\label{def:thresholds}
Under \emph{random ablation} each internal vertex is deleted independently with
probability $q$; the \emph{survival probability} $\Theta(q)$ is the probability that
$S$ remains computable. Under \emph{targeted ablation} an adversary chooses which
vertices to delete; the \emph{targeted critical threshold} $t^*(G)$ is the minimum
number of vertices an adversary must delete to make $S$ non-computable.
\end{definition}

\subsection{Prerequisite results}

We use the following classical theorems without proof.

\begin{theorem}[Chain rule for mutual information {\cite[Thm.~2.5.2]{cover2006elements}}]
\label{thm:chainrule}
For any random variables $S$ and $X=(X_1,X_2)$,
\[
  \II(S;X_1,X_2)=\II(S;X_1)+\II(S;X_2\mid X_1).
\]
Moreover $\II(S;X_2\mid X_1)\ge 0$, with equality if and only if $S\indep X_2\mid X_1$.
\end{theorem}

\begin{theorem}[Fano's inequality {\cite[Thm.~2.10.1]{cover2006elements}}]
\label{thm:fano}
Let $S\to Y\to\widehat{S}$ be a Markov chain with $S$ taking values in a finite set
$\mathcal{S}$, and let $P_e=\Pr[\widehat{S}\neq S]$. Writing $H_2(p)=-p\log_2 p-(1-p)\log_2(1-p)$
for the binary entropy,
\[
  H_2(P_e)+P_e\log_2(\lvert\mathcal{S}\rvert-1)\;\ge\;H(S\mid Y).
\]
\end{theorem}

\begin{theorem}[Menger's theorem, vertex form {\cite[Thm.~3.3.1]{diestel2017graph}}]
\label{thm:menger}
For non-adjacent vertices $\sigma,\tau$ of a graph, the minimum number of vertices
whose deletion separates $\sigma$ from $\tau$ equals the maximum number of internally
vertex-disjoint $\sigma$--$\tau$ paths.
\end{theorem}

We also recall, only as an interpretive cross-check, the partial information
decomposition of \cite{williams2010pid}: for a target $S$ and two sources $R_1,R_2$,
the mutual information $\II(S;R_1,R_2)$ decomposes into nonnegative \emph{redundant},
\emph{unique}, and \emph{synergistic} atoms. We do not use this decomposition in any
proof; the main bounds rest on Shannon information alone.
